% SPDX-License-Identifier: CC-BY-4.0
% Build: make   (needs pdflatex with pgfplots, booktabs, enumitem, microtype)
\documentclass[11pt,a4paper]{article}

\usepackage[T1]{fontenc}
\usepackage{lmodern}
\usepackage{textcomp}
\usepackage{microtype}
\usepackage[margin=2.5cm]{geometry}
\usepackage{booktabs}
\usepackage{tabularx}
\usepackage{array}
\usepackage{enumitem}
\usepackage{xcolor}
\usepackage{pgfplots}
\pgfplotsset{compat=1.18}
\PassOptionsToPackage{hyphens}{url}
\usepackage[hidelinks]{hyperref}
\usepackage{url}
\Urlmuskip=0mu plus 1mu

\setlist{itemsep=2pt,topsep=4pt}
\newcommand{\code}[1]{\texttt{#1}}
\newcolumntype{L}{>{\raggedright\arraybackslash}X}
\renewcommand{\arraystretch}{1.15}

% Generated by analyze.py
\newcommand\fitslope{0.83331}
\newcommand\fitpoints{60}


\title{Why a Capable CPU Boots With Broken Clocks\\[4pt]
  \large The Legion Pro 5 TSC bug and the case for open firmware}
\author{Hasin Ishraq\\
  \small Undergraduate CSE student, Islamic University of Technology (IUT), Bangladesh\\
  \small \href{mailto:hasinishraq001@gmail.com}{\nolinkurl{hasinishraq001@gmail.com}}}
\date{28 September 2026}

\begin{document}
\maketitle

\begin{abstract}
\noindent
On every boot, the firmware of a Lenovo Legion Pro 5 16ADR10 (BIOS RLCN32WW)
hands Linux a processor whose cores disagree about the time by about two
seconds. The system journal recorded this on all 69 unfixed boots with a
complete kernel log between 28 August and 26 September 2026, across six
Linux kernels. The hardware is not at fault: the same cores agree perfectly
after resuming from sleep, and a 0.6-second pre-boot tool that moves one
counter makes Linux accept them on every boot. The bug is a few lines of
firmware logic, but the owner cannot fix it, because the CPU is fused to run
only Lenovo-signed firmware. That is the argument this report makes for open
firmware: source code the community can read and patch, delivered through the
same signed, recoverable update path vendors already use.
\end{abstract}

% ---------------------------------------------------------------------------
\section{The boot history}

The systemd journal kept the kernel log of 80 boots from 28 August to
28 September 2026, and every boot without a fix failed the kernel's TSC
check. Since the pre-boot fix, tscsync~\cite{tscsync}, reached its current
form (v1.3), all 7 boots kept the TSC. Table~\ref{tab:history} groups the
boots by what ran before Linux.

\begin{table}[ht]
\centering
\small
\begin{tabularx}{\textwidth}{@{}l l r L@{}}
\toprule
Phase & Dates (2026) & Boots & What the kernel did \\
\midrule
No fix & 28 Aug -- 26 Sep & 68 & Rejected the TSC on all 67 boots with a full
  log; warp 5.04--6.50 billion cycles \\
tscsync v1.0, measure only & 26 -- 27 Sep & 2 & Rejected; the tool saw all
  31 other cores 5,690,270,7xx cycles ahead of CPU~0, and the kernel measured
  5,690,270,706 \\
tscsync v1.1, sync & 27 Sep & 1 & Rejected: a 25-cycle warp remained \\
tscsync v1.2, sync & 27 Sep & 2 & 1 kept the TSC; 1 rejected it (25-cycle
  warp) \\
tscsync v1.3, sync & 27 Sep & 3 & Kept the TSC after a cold start, a clean
  shutdown and a reboot \\
tscsync v1.4, sync & 27 Sep & 1 & Kept the TSC, also through 2 resumes from
  sleep \\
tscsync v1.5, sync & 27 Sep & 2 & Kept the TSC; 556\,ms per boot \\
tscsync v2.0.0, sync & 28 Sep & 1 & Kept the TSC; all 31 cores within $-15$
  to $+1$ cycles of CPU~0; 579\,ms \\
\bottomrule
\end{tabularx}
\caption{Every boot in the journal, grouped by what ran before Linux. Measure
mode reads the counters and writes nothing, so those boots count as unfixed.}
\label{tab:history}
\end{table}

The one unfixed boot without a warp line has a journal that starts
6.3~seconds in, after the kernel's check. The failures of v1.1 and v1.2 came
from the tool's precision, not the hardware; they set the 10-cycle target
later versions use.

% ---------------------------------------------------------------------------
\section{The fingerprint of a firmware bug}

The warp is exactly proportional to how long the firmware ran, with no fixed
part: 833.31 million cycles per second of firmware time, intercept
$-0.41$\,ms, $R^2 = 0.999999$, largest miss 0.18\,ms
(Figure~\ref{fig:warp}).

\begin{figure}[ht]
\centering
\begin{tikzpicture}
\begin{axis}[
  width=0.9\textwidth, height=7.5cm,
  xmin=0, xmax=8, ymin=0, ymax=8,
  xtick={0,2,4,6,8}, ytick={0,2,4,6,8},
  xlabel={Firmware time before the boot loader (s)},
  ylabel={TSC warp at boot (billion cycles)},
  grid=major, grid style={gray!25},
  legend pos=north west, legend cell align=left,
  tick label style={font=\small}, label style={font=\small},
  legend style={font=\small},
]
\addplot[only marks, mark=*, mark size=1.6pt, color=blue!70!black, fill opacity=0.7]
  table[x=firmware_s, y=warp_gcycles, col sep=comma] {data/warp-vs-firmware.csv};
\addlegendentry{Unfixed boot (\fitpoints{} boots)}
\addplot[domain=0:8, samples=2, dashed, thick, color=red!70!black]
  {\fitslope * x};
\addlegendentry{Line through zero, slope \fitslope{} billion cycles/s}
\end{axis}
\end{tikzpicture}
\caption{TSC warp the kernel measured against systemd's firmware time, for
every unfixed boot that logged both (28 August to 27 September 2026). The 9
unfixed boots without a point are PackageKit offline-update boots, which
reboot before systemd records the firmware time.}
\label{fig:warp}
\end{figure}

The firmware time comes from the firmware's own boot-performance table
(ACPI FPDT), which it measures with CPU~0's TSC. At a TSC rate of about
2.495\,GHz, the warp is one third of CPU~0's own count: when the boot loader
starts, CPU~0's counter reads three quarters of every other core's. A delay
would add a constant; this is a scale error in how CPU~0's counter is set up
or counted.

Three more measurements point at the same place:
\begin{itemize}
\item \textbf{Only CPU~0 is off.} Measured from GRUB before Linux, all 31
  other cores agreed within about 40 cycles, across both chiplets, and all
  were 5.2--5.7 billion cycles ahead of CPU~0.
\item \textbf{Nothing after the firmware changes it.} The pre-boot tool
  measured 5,690,270,7xx cycles; the kernel measured 5,690,270,706 moments
  later on the same boot.
\item \textbf{The resume path is correct.} After a suspend and resume, the
  firmware restarts all counters and every core agrees; the kernel re-checks
  each core and keeps the TSC.
\end{itemize}

% ---------------------------------------------------------------------------
\section{The hardware is capable}

Nothing about the Ryzen 9 8945HX prevents synchronized counters. Four facts
from this machine show it:
\begin{itemize}
\item \textbf{The counters are invariant.} The CPU reports an invariant TSC
  (\code{constant\_tsc}, \code{nonstop\_tsc}): every core's counter runs at a
  fixed rate regardless of power state, so once the counters agree they stay
  in agreement.
\item \textbf{The firmware already gets it right once.} Resuming from sleep
  (S3) runs a different firmware path, which starts all 32 counters together.
  Linux re-checks every core on resume and keeps the TSC.
\item \textbf{One write fixes it.} A 0.6-second pre-boot tool moves only
  CPU~0's counter forward by about 5.2 billion cycles. The kernel then accepts
  the TSC on every boot, and user-space probes see all cores within a few
  dozen cycles for the rest of the session.
\item \textbf{Windows hides the gap.} Windows synchronizes the counters itself
  at boot and exposes that as a boot option, \code{tscsyncpolicy}
  \cite{bcdedit}. A laptop sold and tested with Windows never shows the
  symptom.
\end{itemize}

% ---------------------------------------------------------------------------
\section{Why this BIOS fails and others don't}

On AMD Zen, the firmware alone is responsible for handing over synchronized
counters. This BIOS's cold-boot path leaves CPU~0 behind, and nothing
downstream catches it.

\paragraph{Who is allowed to fix an offset.} Intel CPUs have a per-core
register, \code{IA32\_TSC\_ADJUST}, and Linux uses it to pull a mismatched
core back into line when it comes online~\cite{tscsync-c}. Zen CPUs have no
such register (this machine reports \code{tsc\_adjust=0}). Without it, Linux
tests once and falls back to HPET; the kernel's timekeeping maintainer
rejected writing the counters directly from the kernel as ``flaky and
unreliable''~\cite{lkml}. So on AMD, only firmware can get this right.

\paragraph{What this BIOS does differently.} The data narrows it down, though
not to a line of code: the error is a clean scale factor on CPU~0 alone,
produced during the pre-boot firmware phase and absent from the resume path.
The likeliest explanation is that the cold-boot path initializes or clocks
CPU~0's counter differently from the other 31 cores and never re-synchronizes
them before handing over. Confirming it needs Lenovo's source.

\paragraph{Why it ships anyway.} Windows re-synchronizes the counters itself
at every boot, so a laptop validated on Windows passes. Firmware that syncs
correctly (the resume path here, and most other vendors' boot paths) leaves
nothing for either operating system to repair, which is why most AMD laptops
never show this warning.

\paragraph{It is not the first time.} The same symptom was reported on the
IdeaPad 5 2-in-1 14AHP9 (BIOS P1CN32WW) and on another Legion Pro 5 16ADR10,
with a Ryzen 9 8940HX, on the same RLCN32WW BIOS~\cite{lenovo-forum}. Users
also report that Lenovo fixed the same class of bug by BIOS update on the
ThinkPad A485, E495 and T14s Gen~1 AMD models; those forum pages could not be
opened for this report, so treat that history as reported, not verified.

% ---------------------------------------------------------------------------
\section{What it costs users}

On HPET, every clock read is a trip to a slow timer device, and the kernel
and desktop read the clock tens of thousands of times a second
(Table~\ref{tab:cost}).

\begin{table}[ht]
\centering
\small
\begin{tabularx}{\textwidth}{@{}L L L@{}}
\toprule
Measure on this laptop & HPET (firmware as shipped) & TSC (after the pre-boot
  fix) \\
\midrule
\code{clock\_gettime} from Python, per call & 1,241\,ns & 144\,ns \\
CPU temperature, mostly idle & 66--72\,\textdegree C (owner's observation) &
  59--63\,\textdegree C; 59.9\,\textdegree C measured, 97.6\,\% of the time in
  the deepest idle state \\
\bottomrule
\end{tabularx}
\caption{Cost of the HPET fallback. Both timing figures include Python's own
call overhead, so the raw difference is larger. The temperature comparison is
an owner's observation, not yet a controlled test; tscsync includes a
benchmark script (\code{tools/idle-bench.sh}) for one.}
\label{tab:cost}
\end{table}

Forcing the TSC on anyway (\code{tsc=reliable}) is not a workaround: the
counters really are two seconds apart. The IdeaPad owner who first reported
this found that it causes audio and video stutter during playback and erratic
kernel timers~\cite{lenovo-forum}.

% ---------------------------------------------------------------------------
\section{Why owners cannot fix it themselves}

The fix is small, but on this laptop only Lenovo can ship it: the CPU is
permanently fused to Lenovo's firmware signing key.
\begin{itemize}
\item \textbf{Platform Secure Boot is fused.} The CPU's security processor
  reports \code{fused\_part = 1}. Under AMD Platform Secure Boot, the security
  processor checks the firmware's signature against a key hash burned into the
  chip before any x86 code runs, and the fuses cannot be
  reset~\cite{ioactive,sth}. A patched BIOS would simply not start.
\item \textbf{The flash chip is guarded.} \code{rom\_armor\_enforced = 1}:
  writes to the BIOS chip go through the security processor, so software
  cannot write a modified image. Only an external programmer clipped to the
  board could restore the original.
\item \textbf{The code is closed.} AMD's platform initialization code (AGESA)
  and Lenovo's firmware are proprietary. AMD's open replacement, openSIL, is a
  proof of concept on a server reference board today and does not cover this
  laptop~\cite{opensil}.
\end{itemize}

The security model is sound in itself: it stops malware from implanting
itself into firmware. The problem is that it puts every fix in one vendor's
queue, even when the owner has found the bug, measured it and demonstrated
the fix.

% ---------------------------------------------------------------------------
\section{The case for open firmware that updates safely}

The fix for this bug is probably a few lines in the boot path: the code that
syncs the TSCs after S3 resume already works. Only Lenovo can make that change
today, and owners have reported it on Lenovo's forum and by email. Meanwhile
every owner of this model runs a clock that is wrong on every cold boot. Open
firmware would change who can find and fix bugs like this. It does not have
to make firmware less safe.

\subsection{``Open'' does not mean ``unsigned'' or ``flash it yourself''}

The usual objection is that open firmware means owners flashing random ROMs
and bricking laptops. That mixes up two things:
\begin{itemize}
\item \textbf{Who can read and rebuild the source.} This is what ``open''
  means here.
\item \textbf{Who can sign a build the machine will accept.} This can stay
  with the vendor.
\end{itemize}

AMD Platform Secure Boot and Lenovo's signing keys can stay exactly as they
are. The vendor still builds, signs and ships every update. The difference is
that the source is published and the build is reproducible. Anyone can then
check that the signed binary matches the source, and anyone can find a bug,
write the patch and send it to the vendor. The vendor still decides what
ships. Linux distributions already work this way with signed kernels under
Secure Boot.

\subsection{What ``doesn't brick the device'' needs, and where it already
exists}

Lenovo's own BIOS updates are safe because of the delivery mechanism, not
because the source is closed. Every part of that mechanism already exists in
open form (Table~\ref{tab:safety}).

\begin{table}[ht]
\centering
\small
\begin{tabularx}{\textwidth}{@{}>{\raggedright\arraybackslash}p{3.3cm} L L@{}}
\toprule
Safety property & How it works & Open implementation in use today \\
\midrule
Only vendor-signed images are accepted
  & The update is checked against a key in firmware before it is written
  & Chromebook verified boot checks each firmware copy's signature with a key
    in the read-only boot stub~\cite{chromium}; Dasharo supports verified boot
    and signed releases~\cite{dasharo} \\
\addlinespace
A failed update cannot leave the machine unbootable
  & A read-only recovery region plus two copies of the writable firmware (A/B)
  & Chromebooks: ``The writable portion of the firmware should exist in two
    copies. In the event the first copy is corrupt, the device can boot
    normally off the second copy.''~\cite{chromium} \\
\addlinespace
Updates are applied by the firmware, not by a flashing tool in the OS
  & UEFI capsule update: the OS stages the image and the firmware applies it
    on the next boot
  & fwupd calls the UEFI \code{UpdateCapsule} service~\cite{lvfs}; Dasharo
    ships capsule updates~\cite{dasharo} \\
\addlinespace
Owners get updates without Windows or vendor tools
  & A signed catalog and signed update files
  & The LVFS, used by over one hundred vendors, signs each capsule, and fwupd
    installs it on Linux~\cite{fwupd,lvfs} \\
\addlinespace
Anyone can check the binary against the source
  & Reproducible builds
  & Dasharo documents how to verify signatures and reproduce its
    builds~\cite{dasharo} \\
\bottomrule
\end{tabularx}
\caption{The properties that make a firmware update safe, and open
implementations of each.}
\label{tab:safety}
\end{table}

None of this is experimental:
\begin{itemize}
\item Every Chromebook ships \textbf{coreboot}~\cite{coreboot} with an A/B
  firmware layout and a read-only recovery path, on millions of devices.
\item \textbf{Dasharo} (3mdeb) ships coreboot with EDK~II UEFI on laptops sold
  by NovaCustom and others, with fwupd and capsule updates~\cite{dasharo}.
\item Lenovo already publishes firmware for many ThinkPad models on the
  \textbf{LVFS}. Legion models are generally not there, so Linux owners of
  this laptop cannot even install a BIOS update without Windows or a USB tool.
\end{itemize}

\subsection{What blocks this on AMD laptops today}

The part of the firmware that holds this bug is the CPU and platform
initialization code: AMD's AGESA, integrated by Insyde and Lenovo. That code
is closed.

AMD's open replacement, \textbf{openSIL}, is public and MIT-licensed, but its
repository describes it as a proof of concept for a server reference board
(Genoa), ``NOT for use in production firmware''~\cite{opensil}. It has no
client (laptop) CPU support. Press reports say client support is planned for
a future CPU generation; AMD has not shipped it. Until then, even a vendor
willing to open its firmware could not open the part that syncs the TSCs.

So the realistic path has two steps:
\begin{enumerate}
\item \textbf{Now:} vendors publish the platform code they own, ship through
  the LVFS, and test cold boot on Linux.
\item \textbf{When openSIL covers client CPUs:} laptops built on coreboot or
  EDK~II with openSIL, vendor-signed, with A/B recovery, delivered as
  capsules. A bug like this one could then be found, fixed and reviewed in
  public, and still reach owners only through a signed, recoverable update.
\end{enumerate}

\subsection{Why it matters beyond one laptop}

This bug shows what closed firmware costs:
\begin{itemize}
\item \textbf{The hardware is fine and the fix is known.} The same firmware
  already syncs the TSCs correctly on resume.
\item \textbf{Windows hides the problem.} Windows resyncs the TSCs itself, so
  a vendor that tests only on Windows never sees it.
\item \textbf{Only the vendor can fix it, and the vendor has little reason
  to.} Owners are left with a pre-boot workaround (tscsync) that has to be
  maintained, signed and installed outside the firmware.
\end{itemize}

With open source, reproducible builds and vendor signing, the person who
found the bug could have sent the fix, and the vendor could have reviewed,
signed and shipped it through the same safe update path it uses today.

% ---------------------------------------------------------------------------
\section{Recommendations}

\subsection{For Lenovo and other vendors}
\begin{enumerate}
\item \textbf{Fix the cold-boot TSC sync.} Apply the same sync between CPU~0
  and the other cores that the S3 resume path already does, before handing
  off to the operating system. The kernel's check
  (\code{check\_tsc\_sync\_source}) should pass on every core with no help
  from the OS.
\item \textbf{Test cold boot on Linux.} A single boot of a current Fedora or
  Ubuntu live image shows \code{TSC warp} in the kernel log. Windows hides the
  bug because it resyncs the TSCs itself.
\item \textbf{Ship firmware through the LVFS} for consumer lines too, not
  only ThinkPads, so Linux owners get fixes through fwupd.
\item \textbf{Publish the platform code you own}, with reproducible builds,
  and move to openSIL when AMD supports client CPUs.
\end{enumerate}

\subsection{For owners of affected machines}
\begin{itemize}
\item \textbf{Check} the kernel log of the current boot:\\
  \code{journalctl -k -b | grep -iE \textquotesingle tsc|clocksource\textquotesingle}\\
  \code{Marking TSC unstable} or \code{TSC warp}, followed by
  \code{Switched to clocksource hpet}, means you are affected.
\item \textbf{Report it} to the vendor with the model, BIOS version and that
  log line. A workaround does not replace a firmware fix.
\item \textbf{Work around it} with tscsync~\cite{tscsync}: install it in
  measure mode first, then switch to sync mode. It writes nothing to firmware
  and disables itself if a boot fails.
\item \textbf{After every BIOS update}, switch back to measure mode for one
  cold boot. If every core shows \code{OK}, the vendor has fixed the bug and
  tscsync can be removed.
\end{itemize}

% ---------------------------------------------------------------------------
\appendix
\section{Methodology and data}

All data and scripts are in the tscsync repository under
\code{docs/report/}~\cite{tscsync}.
\begin{itemize}
\item \textbf{Boot history} (\code{data/boots.csv}, written by
  \code{extract-boots.sh}): one row per boot in the systemd journal of this
  machine, every boot since installation. For each boot: the kernel version,
  the warp from \code{Measured N cycles TSC warp between CPUs}, systemd's
  firmware time from \code{Startup finished in \dots{} (firmware)} (which
  systemd reads from the ACPI FPDT), the tscsync version and mode, and
  whether the kernel kept the TSC.
\item \textbf{Warp against firmware time} (\code{analyze.py}): an ordinary
  least-squares fit of the warp against the firmware time over every unfixed
  boot that logged both (\fitpoints{} boots). It writes the plotted points to
  \code{data/warp-vs-firmware.csv}.
\item \textbf{Per-core offsets:} measured before Linux starts by tscsync in
  measure mode, and after boot by \code{tools/tscprobe.c}. Both use a
  paired-read test between CPU~0 and each other core, giving the offset in
  cycles and its direction.
\item \textbf{Resume path:} the kernel's TSC check on S3 resume, with and
  without the workaround.
\item \textbf{Limits:} one machine and one BIOS version (RLCN32WW). The other
  models are listed from owner reports and were not measured here.
\end{itemize}

% ---------------------------------------------------------------------------
\begin{thebibliography}{99}
\small

\bibitem{tscsync}
H.~Ishraq, \emph{tscsync: fix TSC warp between CPUs before Linux boots},
\url{https://github.com/Zanasin/tscsync} (see \code{docs/DESIGN.md} and
\code{docs/HARDWARE.md}).

\bibitem{tscsync-c}
Linux kernel, \code{arch/x86/kernel/tsc\_sync.c},
\url{https://github.com/torvalds/linux/blob/master/arch/x86/kernel/tsc_sync.c}.

\bibitem{lkml}
Linux kernel mailing list, discussion of TSC synchronization on AMD CPUs
without \code{TSC\_ADJUST}, August 2022,
\url{https://lkml.iu.edu/hypermail/linux/kernel/2208.3/00630.html}.

\bibitem{bcdedit}
Microsoft, \emph{BCDEdit /set} (\code{tscsyncpolicy}),
\url{https://learn.microsoft.com/en-us/windows-hardware/drivers/devtest/bcdedit--set}.

\bibitem{lenovo-forum}
Lenovo Community, \emph{1 second TSC warp on IdeaPad 5 2-in-1 14AHP9 due to
bad initialization in BIOS forces HPET fallback},
\url{https://forums.lenovo.com/t5/Other-Linux-Discussions/1-second-TSC-warp-on-IdeaPad-5-2-in-1-14AHP9-due-to-bad-initialization-in-BIOS-forces-HPET-fallback/m-p/10021248}.

\bibitem{ioactive}
IOActive, \emph{Exploring AMD Platform Secure Boot},
\url{https://www.ioactive.com/exploring-amd-platform-secure-boot/}.

\bibitem{sth}
ServeTheHome, \emph{AMD PSB vendor locks EPYC CPUs for enhanced security at a
cost},
\url{https://www.servethehome.com/amd-psb-vendor-locks-epyc-cpus-for-enhanced-security-at-a-cost/}.

\bibitem{opensil}
AMD, \emph{openSIL}, \url{https://github.com/openSIL/openSIL}.

\bibitem{coreboot}
coreboot documentation, \url{https://doc.coreboot.org/}.

\bibitem{dasharo}
3mdeb, \emph{Dasharo documentation}, \url{https://docs.dasharo.com/}.

\bibitem{chromium}
The Chromium Projects, \emph{Firmware boot and recovery},
\url{https://www.chromium.org/chromium-os/chromiumos-design-docs/firmware-boot-and-recovery/}.

\bibitem{fwupd}
fwupd, \url{https://fwupd.org/}.

\bibitem{lvfs}
Linux Vendor Firmware Service, \emph{Introduction},
\url{https://lvfs.readthedocs.io/en/latest/intro.html}.

\end{thebibliography}

\end{document}
